\chapter{验收与发布}

\section{自动化测试}

\subsection{静态代码验证}

在动态测试之前，团队采用了两类静态验证手段。

\textbf{TypeScript 类型检查}：前端生产构建将静态类型检查作为必要前置步骤，任何类型不兼容错误均会阻断构建。该检查覆盖页面组件、状态管理、通信封装和数据类型定义，用于验证前端数据结构与后端接口契约的一致性。

\textbf{代码走查}：每个开发迭代完成后，团队成员对新增实现进行交叉检查，重点核查身份认证、资源归属、多用户数据隔离和分类异常处理，发现的问题在当前迭代内修正。

\subsection{动态测试}

系统后端建立了覆盖认证、导入、交易、分类、解析、分析、人格、组织、报表和故障重试等领域的自动化测试体系。表~\ref{tab:test-files} 按功能域汇总测试内容与数量，避免以源文件清单代替测试说明。

\begin{longtable}{p{4.2cm} p{7.2cm} p{2.1cm}}
\caption{自动化测试覆盖范围}
\label{tab:test-files}\\
\toprule
\textbf{功能域} & \textbf{主要验证内容} & \textbf{测试数} \\
\midrule
\endfirsthead
\multicolumn{3}{c}{\small 表~\thetable\quad 自动化测试覆盖范围（续）} \\
\toprule
\textbf{功能域} & \textbf{主要验证内容} & \textbf{测试数} \\
\midrule
\endhead
\midrule
\multicolumn{3}{r}{\small 接下页} \\
\endfoot
\bottomrule
\endlastfoot
认证与数据隔离 & 密码摘要、身份凭证、登录流程和受保护资源访问 & 5 \\
\hline
导入、解析与去重 & 多平台账单识别、文本规范化、重复交易识别、批次进度和部分失败处理 & 10 \\
\hline
分类与故障重试 & 提示信息、分类缓存、规则降级、超时重试、失败恢复和可靠性指标 & 14 \\
\hline
交易与报表 & 交易筛选、分页、人工分类、跨用户保护、报表生成与下载 & 7 \\
\hline
数据分析与消费人格 & 财务聚合、日期与分类筛选、人格维度、问卷偏差及边界条件 & 10 \\
\hline
家庭组织 & 组织创建、成员权限、家庭聚合、角色区分和个人模式兼容 & 13 \\
\end{longtable}

\section{需求—设计—测试追踪}

表~\ref{tab:fr-traceability} 和表~\ref{tab:nfr-traceability} 将需求与设计模块、测试证据对应起来。“自动化通过”表示存在可重复执行的测试；“手工验收”表示通过第 7.4 节的用户流程检查；“设计目标”表示尚未完成专项测量，不能解释为已经达到。

\begingroup
\footnotesize
\setlength{\tabcolsep}{3pt}
\begin{longtable}{@{}p{1.0cm} p{2.8cm} p{3.8cm} p{3.6cm} p{1.3cm}@{}}
\caption{功能需求追踪矩阵}
\label{tab:fr-traceability}\\
\toprule
\textbf{ID} & \textbf{设计/实现} & \textbf{测试证据} & \textbf{验收场景} & \textbf{状态} \\
\midrule
\endfirsthead
\multicolumn{5}{c}{\small 表~\thetable\quad 功能需求追踪矩阵（续）} \\
\toprule
\textbf{ID} & \textbf{设计/实现} & \textbf{测试证据} & \textbf{验收场景} & \textbf{状态} \\
\midrule
\endhead
\midrule
\multicolumn{5}{r}{\small 接下页} \\
\endfoot
\bottomrule
\endlastfoot
FR-01 & 身份认证与前端登录状态管理 & 注册、登录、凭证校验与页面访问控制测试 & 步骤 1 & 自动化通过 \\
\hline
FR-02 & 当前用户识别与资源归属过滤 & 交易、导入、报表越权测试；组织权限测试 & 步骤 1、5、9、12 & 自动化通过 \\
\hline
FR-03 & 账单导入、批次与文件管理 & 导入流程、批次进度和用户隔离测试 & 步骤 2、3 & 自动化通过 \\
\hline
FR-04 & 平台识别、账单解析与文本规范化 & 微信、支付宝样例及规范化稳定性测试 & 步骤 2、3 & 自动化通过 \\
\hline
FR-05 & 复合分类与分类历史 & 缓存、模型、规则降级和结果记录测试 & 步骤 6 & 自动化通过 \\
\hline
FR-06 & 交易查询、筛选与分页 & 条件过滤、分页、身份校验和数据隔离测试 & 步骤 5、6 & 自动化通过 \\
\hline
FR-07 & 人工分类与校正工作台 & 分类更新、复核状态和越权访问测试 & 步骤 7 & 自动化通过 \\
\hline
FR-08 & 系统分类与自定义分类管理 & 当前以页面操作和接口联调验证 & 步骤 8 & 手工验收 \\
\hline
FR-09 & 财务聚合分析与看板 & 日期、分类筛选及汇总指标测试 & 步骤 4 & 自动化通过 \\
\hline
FR-10 & 消费人格与问卷分析 & 人格维度、问卷偏差和边界条件测试 & 步骤 10、11 & 自动化通过 \\
\hline
FR-11 & 家庭组织与基于角色的权限控制 & 组织创建、成员管理和角色权限测试 & 步骤 12 & 自动化通过 \\
\hline
FR-12 & 成员 ID 范围解析、家庭聚合参数 & 组织聚合与个人模式兼容测试 & 步骤 13 & 自动化通过 \\
\hline
FR-13 & 报表构建、记录与下载 & 报表筛选、生成和跨用户下载保护测试 & 步骤 9 & 自动化通过 \\
\hline
FR-14 & 用户分类设置与可替换分类策略 & 分类策略测试、前端类型检查和设置页验收 & 步骤 14 & 自动化 + 验收 \\
\hline
FR-15 & 重试任务处理器与管理员恢复功能 & 重试状态、进度、失败恢复和超时测试 & 步骤 14 & 自动化通过 \\
\end{longtable}
\endgroup

\begingroup
\footnotesize
\setlength{\tabcolsep}{3pt}
\begin{longtable}{@{}p{1.0cm} p{2.9cm} p{4.0cm} p{3.2cm} p{1.3cm}@{}}
\caption{非功能需求追踪矩阵}
\label{tab:nfr-traceability}\\
\toprule
\textbf{ID} & \textbf{设计机制} & \textbf{验证方式} & \textbf{当前结论} & \textbf{状态} \\
\midrule
\endfirsthead
\multicolumn{5}{c}{\small 表~\thetable\quad 非功能需求追踪矩阵（续）} \\
\toprule
\textbf{ID} & \textbf{设计机制} & \textbf{验证方式} & \textbf{当前结论} & \textbf{状态} \\
\midrule
\endhead
\midrule
\multicolumn{5}{r}{\small 接下页} \\
\endfoot
\bottomrule
\endlastfoot
NFR-01 & 密码哈希、JWT、RBAC、资源归属 & 安全、认证、交易与组织 API 测试 & 已覆盖身份与普通越权；上传限制和速率限制待补 & 部分验证 \\
\hline
NFR-02 & 用户 ID 隔离、家庭成员显式聚合 & 导入、交易、报表和组织隔离测试 & 核心私有资源具备隔离证据 & 自动化通过 \\
\hline
NFR-03 & 完整用户流程与个人/家庭双口径 & 14 步手工验收与组织兼容测试 & 主要任务可完成 & 自动化 + 验收 \\
\hline
NFR-04 & 超时、缓存、日期过滤和进度反馈 & 分类器超时测试；尚无标准化性能基准 & 超时机制已验证，P95 和千笔解析指标仍是目标 & 设计目标 \\
\hline
NFR-05 & 本地、Docker Compose、裸机文档 & 配置审查、前端构建、部署手册 & 具备三种运行说明；未执行持续部署测试 & 文档验证 \\
\hline
NFR-06 & 路由、服务、模型与 Schema 分层 & 代码走查、模块和接口文档 & 主要职责边界清晰 & 代码审查 \\
\hline
NFR-07 & 后端自动化测试、前端单元测试与类型检查 & 本章测试结果与生产构建验证 & 核心后端与前端基础能力具备回归测试 & 自动化通过 \\
\hline
NFR-08 & 可替换解析、分类、报表构建和异步任务边界 & 解析、分类和报表测试 & 主要变化点已有扩展边界，异步报表处理尚未完成 & 部分验证 \\
\end{longtable}
\endgroup


\section{测试结果}

本轮验证于 2026 年 6 月 15 日在 macOS、Python 3.11.14 和 Node.js 25.2.1 环境执行。全部 59 个后端测试用例通过，后端整体覆盖率为 75\%。测试覆盖了各主要业务模块的正常流程、权限边界和异常处理。

前端除静态类型检查外，还建立了单元测试体系，覆盖认证状态管理、登录与登出、通信层身份凭证附加、页面访问控制、组织数据请求以及重试进度解析和异常恢复。全部 26 个前端测试用例通过。

前端构建验证包括 TypeScript 类型检查与 Vite 生产构建，两个阶段均通过，无类型错误或编译错误。构建结果显示部分前端资源包超过 500 KB，主要来源于 Element Plus 与 ECharts，属于后续性能优化项，不影响功能完整性。依赖审计同时报告 8 个已知漏洞（2 个中危、6 个高危）；本次未在缺少回归评审的情况下自动升级依赖，而是将其列为发布前必须复核的安全项。

\begin{table}[H]
\centering
\caption{测试结果汇总}
\label{tab:test-results}
\begin{tabular}{p{6cm} p{3cm} p{4.5cm}}
\toprule
\textbf{测试类型} & \textbf{结果} & \textbf{说明} \\
\midrule
后端单元与集成测试 & 59/59 通过 & 覆盖六个主要功能域 \\
\hline
后端代码覆盖率 & 75\% & 覆盖率工具统计 \\
\hline
前端单元测试 & 26/26 通过 & 覆盖认证、通信、导航和重试进度 \\
\hline
前端 TypeScript 类型检查 & 通过 & 未发现类型错误 \\
\hline
前端 Vite 生产构建 & 通过 & 静态资源构建成功 \\
\hline
前端资源包体积优化 & 待优化 & Element Plus 与 ECharts 相关资源超过 500KB \\
\hline
前端依赖审计 & 8 项待复核 & 2 个中危、6 个高危，需评估升级影响 \\
\bottomrule
\end{tabular}
\end{table}

\section{验收测试}

验收测试以手工联调方式执行，覆盖从注册到报表生成的完整用户流程，共 14 个验收步骤。

\begin{longtable}{p{0.8cm} p{5cm} p{7.5cm}}
\caption{验收测试执行清单}
\label{tab:acceptance-tests}\\
\toprule
\textbf{步骤} & \textbf{操作} & \textbf{预期结果} \\
\midrule
\endfirsthead
\multicolumn{3}{c}{\small 表~\thetable\quad 验收测试执行清单（续）} \\
\toprule
\textbf{步骤} & \textbf{操作} & \textbf{预期结果} \\
\midrule
\endhead
\midrule
\multicolumn{3}{r}{\small 接下页} \\
\endfoot
\bottomrule
\endlastfoot
1 & 注册新用户并登录 & 注册成功，进入财务看板 \\
\hline
2 & 进入导入页面，上传微信或支付宝账单 & 文件上传成功，显示导入批次进度 \\
\hline
3 & 检查最近导入结果与批次进度 & 批次由处理中变为完成，交易计数正确 \\
\hline
4 & 进入财务看板查看概览 & 收入、支出、储蓄率、分类占比与主要商户正常显示 \\
\hline
5 & 进入交易列表，按条件筛选 & 筛选结果正确，分页可用 \\
\hline
6 & 打开交易详情抽屉 & 显示分类来源、置信度与分类理由 \\
\hline
7 & 进入校正工作台，确认一笔待校正交易 & 交易从待校正队列移除，最终分类已更新 \\
\hline
8 & 新增一个用户自定义分类 & 分类列表中显示新分类，校正工作台中可选 \\
\hline
9 & 进入报表中心，生成 PDF & 生成成功，文件可下载，内容与交易数据一致 \\
\hline
10 & 进入消费人格页面 & 人格画像卡片、四维雷达图、财务健康评分正常渲染 \\
\hline
11 & 完成人格问卷自评 & 显示主观认知与交易数据画像的偏差对比 \\
\hline
12 & 创建家庭组织并添加成员 & 组织创建成功，成员列表中显示被添加用户 \\
\hline
13 & 查看家庭组织财务看板 & 显示家庭聚合收入、支出、交易数与家庭人格雷达图 \\
\hline
14 & 查看系统设置页面 & 管理员可查看并处理重试任务，普通用户仅可查看本人配置 \\
\end{longtable}

每个验收步骤均在桌面视口（1440x900）下完成验证，确保页面无中文乱码、布局正常、交互响应正确。窄屏视口（390x844）下验证导航栏与内容区无严重重叠，表格与图表可正常展示。

\section{发布状态}

系统具备可复现的 Docker Compose 与 Nginx 课程部署方案，各组件职责如下：

\begin{itemize}
  \item \textbf{前端静态资源}：Vue 3 应用通过 Vite 构建为静态文件，由 Nginx 直接托管提供服务。
  \item \textbf{Nginx 反向代理}：提供静态文件服务，并将业务请求转发至后端应用。
  \item \textbf{FastAPI 后端}：提供 REST API，承担身份校验、业务处理和数据访问。
  \item \textbf{PostgreSQL}：持久化存储用户、交易、分类、组织、报表等业务数据。
  \item \textbf{Redis}：为可选的异步任务处理提供消息传递与结果存储。当前主要业务仍可采用应用内任务或同步执行。
  \item \textbf{重试任务处理器}：在应用运行期间依次处理等待重试的外部模型分类请求。
\end{itemize}

管理员账号通过受控方式创建，不提供公开注册入口。数据库连接、模型服务密钥和 Redis 地址等敏感配置由部署环境统一管理，不写入源代码或公开配置。

系统当前版本已交付的完整功能包括：用户认证与数据隔离、微信和支付宝账单导入与解析、大模型语义分类与人工校正、财务看板、消费人格画像与财务健康评分、家庭组织管理与聚合分析、PDF 报表生成与下载、多种分类服务切换与失败重试。系统具备可复现的部署基础，但速率限制、自动备份、监控告警和报表异步处理流程仍属于后续完善项。
